\documentclass{article}
% ---------- Codificación e idioma ----------
\usepackage[utf8]{inputenc}
\usepackage[T1]{fontenc}
\usepackage{csquotes} 

% ---------- Matemáticas ----------
\usepackage{amsmath, amssymb} 

% ---------- Tipografía ----------
\usepackage[ttdefault=true]{AnonymousPro}
\usepackage{parskip}
\usepackage{microtype}

% ---------- Geometría ----------
\usepackage[margin=2.5cm]{geometry}

% ---------- Listas ----------
\usepackage{enumitem}
\setlist[itemize]{leftmargin=*, itemsep=2pt, topsep=3pt}
\setlist[enumerate]{leftmargin=*, itemsep=2pt, topsep=3pt}

% ---------- Cajas de color ----------
\usepackage{tcolorbox}
\usepackage{xcolor}

% ---------- Gráficos ----------
\usepackage{graphicx}
\usepackage{float}

% ---------- Bibliografía ----------
\usepackage[backend=biber, style=ieee, sorting=none]{biblatex}
\addbibresource{references.bib}

% ---------- Hipervínculos (cargar al final) ----------
\usepackage[hidelinks]{hyperref}

% COLORES INSTITUCIONALES
\definecolor{unamazul}{RGB}{0, 83, 155}
\definecolor{unagris}{RGB}{240, 240, 240}

% COMANDOS PERSONALIZADOS
\newcommand{\captura}[1]{
    \begin{center}
        \includegraphics[width=0.7\textwidth]{imgs/#1}
    \end{center}
}

\newcommand{\codigo}[1]{\texttt{#1}}

% PORTADA
\newcommand{\portadapractica}[5]{
\begin{titlepage}
    \centering
    \vspace*{1cm}
    {\scshape\LARGE Universidad Nacional Autónoma de México\par}
    \vspace{0.4cm}
    {\scshape\large Facultad de Ingeniería\par}
    \vspace{2.5cm}
    {\color{unamazul}\rule{\textwidth}{1.5pt}\par}
    \vspace{0.8cm}
    {\Huge\bfseries #2\par}
    \vspace{0.8cm}
    {\color{unamazul}\rule{\textwidth}{1.5pt}\par}
    \vspace{1.5cm}
    {\Large\itshape #1\par}
    \vspace{1.2cm}
    {\large\textbf{Profesor:} #3\par}
    \vspace{2.5cm}
    {\large\bfseries Integrantes\par}
    \vspace{0.6cm}
    {\large #4\par}
    \vfill
    {\large #5\par}
    \vspace{1cm}
\end{titlepage}
}

\begin{document}

\portadapractica
    {Programación Orientada a Objetos}
    {Investigacion del patrón de diseño: FLyweight}
    {Edgar Tista García}
    {Robles Cancino Arturo Andres\\
    Rodríguez Cárcamo Ditter Emmanuel\\
    Santamaría Viveros César Yurem}
    {2 de Noviembre, 2026}

\section{Casos de uso}
Segun el autor de Design Patterns Elements of Reusable Object-Oriented Software se dice que se debería emplear cuadno se cumplen las sigueintes caracteristicas:
\begin{enumerate}
    \item La aplicación usa una gran cantidad de objetos.
    \item Los costos de almacenamiento son altos debido a la cantidad de objetos.
    \item La mayor parte del estado de los objetos puede hacerse extrínseco.
    \item Muchos grupos de objetos pueden reemplazarse por relativamente pocos objetos compartidos una vez que se elimina el estado extrínseco.
    \item La aplicación no depende de la identidad de los objetos. Dado que los flyweights pueden compartirse, las pruebas de identidad devolverán verdadero para objetos conceptualmente distintos.
\end{enumerate}
El libro lo advierte antes de listar los criterios. No basta con que uno o dos se cumplan, sino que el patrón solo se justifica cuando los cinco se sostienen por la cantidad de abstraccióny división de entidades del software.

\section*{Referencias}
\nocite{*}
\printbibliography[heading=none]

\end{document}